\section{Objectifs}
\begin{itemize}
\item Utiliser un vecteur normal à une droite.
\item Déterminer un projeté orthogonal.
\item Reconnaître et exploiter l’équation d’un cercle.
\end{itemize}

\section{Prérequis}
Coordonnées, équations de droites, produit scalaire.

\section{Activité de découverte}
Une équation cartésienne de droite peut être lue comme une condition d’orthogonalité avec un vecteur normal.

\section{Vocabulaire}
Vecteur normal, équation cartésienne, projection orthogonale, cercle, centre, rayon.

\section{Cours}
Dans un repère orthonormé, le vecteur $(a,b)$ est normal à la droite d’équation
\[
ax+by+c=0.
\]

Si la droite passe par $A(x_A,y_A)$, une équation est
\[
a(x-x_A)+b(y-y_A)=0.
\]

Le projeté orthogonal $H$ de $P$ sur une droite $d$ vérifie $H\in d$ et $\overrightarrow{PH}$ est colinéaire à un vecteur normal à $d$.

Le cercle de centre $C(a,b)$ et de rayon $R$ a pour équation
\[
(x-a)^2+(y-b)^2=R^2.
\]

\section{Exemples corrigés}
La droite $2x-3y+5=0$ admet $(2,-3)$ comme vecteur normal.
L’équation $x^2+y^2-4x+6y-12=0$ devient $(x-2)^2+(y+3)^2=25$.

\section{Méthode}
Identifier d’abord l’objet géométrique puis traduire ses propriétés en équations.

\section{Erreurs fréquentes}
\begin{itemize}
\item Prendre $(a,b)$ pour un vecteur directeur de $ax+by+c=0$.
\item Oublier l’hypothèse de repère orthonormé.
\item Lire un rayon directement dans une équation développée.
\end{itemize}

\section{Synthèse}
Vecteur normal, orthogonalité et équations cartésiennes permettent de traiter efficacement droites, projections et cercles.
